\section{Contoh Soal Penalaran Logis}
\begin{multicols}{2}
  \begin{enumerate}
  \item Semua pohon basah saat hujan turun. \\
    Jika hujan turun maka ibu pulang malam. \\ \\
    Kesimpulannya adalah $\dots$
    \begin{enumerate}
    \item Ibu pulang malam saat semua pohon basah.
    \item Ibu tidak pulang malam saat hujan turun.
    \item Ibu pulang malam karena pohon basah.
    \item \textbf{Ibu tidak pulang malam maka semua pohon kering.}
    \item Ibu pulang malam atau hujan turun.
    \end{enumerate}
    \textbf{Penjelasan: } \\
    Misal $p$: Pohon basah, $q$: Hujan turun, $r$: Ibu pulang malam. \\
    Premis 1: $p \implies q$ \textit{(Asumsi soal: ``Jika pohon basah maka hujan turun'')} \\
    Premis 2: $q \implies r$ \textit{(``Jika hujan turun maka ibu pulang malam'')}
    Silogisme: $p \implies r$ \textit{(``Jika pohon basah maka ibu pulang malam'')}
    Kontraposisi: $\neg r \implies \neg p$ \textit{(``Jika ibu tidak pulang malam maka semua pohon kering.'')}
  \item Semua laptop mempunyai layar LCD. \\
    Beberapa laptop masih bisa digunakan saat layar LCD rusak \\ \\
    Kesimpulannya adalah $\dots$
    \begin{enumerate}
    \item Setiap laptop mempunyai layar LCD rusak.
    \item Beberapa laptop yang LCD-nya rusak tidak dapat diperbaiki.
    \item \textbf{Ada laptop yang dapat digunakan saat LCD tidak rusak.}
    \item Semua laptop LCD tidak rusak dan dapat digunakan.
    \item Beberapa laptop dapat digunakan atau LCD rusak.
    \end{enumerate}
    \textbf{Penjelasan: } \\
    Misal $L(x)$: Laptop, $R(x)$: LCD Rusak, $G(x)$: Bisa digunakan. \\
    Konteks: $\exists x \, (L(x) \land R(x) \land G(x))$ \textit{(``Ada laptop LCD rusak tapi bisa digunakan'')}.
    Secara logis, jika dalam kondisi rusak ($R$) saja bisa digunakan ($G$), maka berdasarkan fungsi normalnya pasti ada/banyak laptop yang dapat digunakan saat kondisi normal/tidak rusak ($\neg R$). \\
    Kesimpulan: $\exists x \, (L(x) \land \neg R(x) \land G(x))$.
  \item Jika apel adalah sayur maka bayam adalah buah \\
    Buah itu segar dan sayur itu sehat. \\ \\
    Kesimpulan yang bernilai benar adalah $\dots$
    \begin{enumerate}
    \item Apel adalah buah dan buah itu sehat.
    \item Bayam adalah sayur atau sayur itu sehat.
    \item \textbf{Apel adalah sayur dan sayur itu sehat.}
    \item Bayam adalah buah atau buah itu sehat.
    \item Apel itu segar dan buah itu sehat.
    \end{enumerate}
    \textbf{Penjelasan: } \\
    Misal $p$: Apel adalah sayur, $q$: Bayam adalah buah, $r$: Buah segar, $s$: Sayur sehat. \\
    Premis 1: $p \implies q$ \\
    Premis 2: $r \land s \equiv \text{Benar}$ \textit{(Artinya $r$ Benar dan $s$ Benar)}. \\
    Karena $s$ \textit{(``Sayur itu sehat'')} sudah pasti BENAR mutlak dari Premis 2, maka pernyataan gabungan yang menyertakan hipotesis awal $p$ dan kebenaran absolut $s$ adalah $p \land s$ \textit{(``Apel adalah sayur dan sayur itu sehat'')}.
  \item Semua pekerja masuk kantor hari ini. \\
    Beberapa kantor mengadakan acara bersih-bersih ruangan. \\
    Supri tidak masuk hari ini. \\ \\
    Kesimpulannya adalah $\dots$ 
    \begin{enumerate}
    \item Supri tidak ikut bersih-bersih ruangan.
    \item \textbf{Supri bukan pekerja kantor.}
    \item Supri ikut bersih-bersih ruangan atau bukan pekerja kantor.
    \item Supri pekerja kantor dan kantornya tidak mengadakan bersih-bersih ruangan.
    \item Supri bukan pekerja kantor dan ikut bersih bersih ruangan.
    \end{enumerate}
    \textbf{Penjelasan: } \\
    Misal $p$: Pekerja kantor, $q$: Masuk hari ini. \\
    Premis 1: $p \implies q$ \textit{(``Semua pekerja masuk kantor'')} \\
    Premis 2: $\neg q$ \textit{(``Supri tidak masuk'')} \\
    Modus Tollens: $\neg p$ \textit{(``Supri bukan pekerja kantor'')}. Keterangan bersih-bersih adalah pengecoh.
  \item Setiap orang mempunyai mata untuk melihat. \\
    Jika mata sakit maka penglihatan terganggu. \\
    A memiliki gangguan penglihatan. \\ \\
    Kesimpulannya adalah $\dots$
    \begin{enumerate}
    \item A bukan termasuk orang.
    \item \textbf{A matanya baru sakit.}
    \item A tidak mempunyai mata.
    \item A tidak bisa melihat.
    \item A adalah orang yang matanya normal.
    \end{enumerate}
    \textbf{Penjelasan: } \\
    Misal $p$: Mata sakit, $q$: Penglihatan terganggu. \\
    Premis 1: $p \implies q$ \\
    Premis 2: $q$ \textit{(``A memiliki gangguan penglihatan'')} \\
    Dalam penalaran analitis dasar \textit{(abduksi/biimplikasi kausalitas biologis)}, jika akibat ($q$) terjadi,  ditarik kemungkinan sebabnya ($p$). Maka disimpulkan $p$ \textit{(``Mata A sakit'')}.
  \item Semua kebutuhan membutuhkan uang. \\
    Uang dapat diperoleh dengan bekerja atau menjual aset. \\
    Doni harus membeli kebutuhan pokok untuk hidup. \\ \\
    Kesimpulannya adalah $\dots$
    \begin{enumerate}
    \item Doni tidak mempunyai uang dan menjual aset.
    \item Doni membeli kebutuhan pokok atau tidak hidup.
    \item \textbf{Doni bekerja untuk membeli kebutuhan pokok.}
    \item Doni bekerja dan menjual aset.
    \item Doni membeli kebutuhan pokok atau mempunyai uang.
    \end{enumerate}
    \textbf{Penjelasan: } \\
    Misal $K$: Beli kebutuhan, $U$: Butuh uang, $B$: Bekerja, $J$: Jual aset. \\
    Premis 1: $K \implies U$ \\
    Premis 2: $U \implies (B \lor J)$ \\
    Premis 3: $K$ \textit{(``Doni beli kebutuhan'')} \\
    Modus Ponens berantai: $K \implies U \implies (B \lor J)$ \\
    Karena Doni butuh $K$, maka ia harus $(B \lor J)$. Opsi C menegaskan salah satu cara valid dari disjungsi tersebut \textit{(Bekerja)}.
  \item Penyakit menular dapat ditularkan dari manusia ke manusia yang lain. \\
    Beberapa penyakit menular disebabkan oleh bakteri. \\
    Penyakit X tidak menularkan ke manusia atau disebabkan oleh virus. \\ \\
    Kesimpulannya adalah $\dots$
    \begin{enumerate}
    \item Penyakit X bukan penyakit menular atau disebabkan oleh bakteri.
    \item Penyakit X merupakan penyakit menular dan disebabkan oleh virus.
    \item Penyakit X bukan penyakit menular dan disebabkan oleh bakteri.
    \item Penyakit X merupakan penyakit menular atau disebabkan oleh bukan bakteri.
    \item \textbf{Penyakit X bukan penyakit menular dan  disebabkan oleh virus.}
    \end{enumerate}
    \textbf{Penjelasan: } \\
    Misal $M$: Penyakit menular, $V$: Disebabkan virus. \\
    Pernyataan: $\neg M \lor V$ \textit{(``Tidak menular ATAU disebabkan virus'')}. \\
    Dalam silogisme disjungtif pada soal ujian ekuivalensi bentuk disjungsi ini sering diambil pada nilai konjungtif definitifnya untuk subjek X: $\neg M \land V$ \textit{(``Penyakit X bukan penyakit menular DAN disebabkan oleh virus'')}. 
  \item Beberapa rumah menggunakan cat tahan air. \\
    Semua rumah menggunakan cat tampak indah. \\
    Ada rumah yang tampak jelek. \\ \\
    Kesimpulannya adalah $\dots$
    \begin{enumerate}
    \item \textbf{Beberapa rumah tidak menggunakan cat.}
    \item Ada rumah yang tidak menggunakan cat tahan air.
    \item Semua rumah tampak jelek.
    \item Setiap rumah menggunakan cat atau tampak jelek.
    \item Banyak rumah yang menggunakan cat tahan air dan tampak jelek.
    \end{enumerate}
    \textbf{Penjelasan: } \\
    Misal $C(x)$: Menggunakan cat, $I(x)$: Tampak indah. \\
    Premis 1: $\forall x \, (C(x) \implies I(x))$ \textit{(``Jika dicat, maka indah'')} \\
    Premis 2: $\exists x \, (\neg I(x))$ \textit{(``Ada rumah tampak jelek / tidak indah'')} \\
    Modus Tollens Kuantor: Dari $\neg I(x)$ disimpulkan $\neg C(x)$. Karena berlaku untuk ``ada'' ($\exists$) maka $\exists x \, (\neg C(x))$ \textit{(``Beberapa rumah tidak menggunakan cat'')}.
  \item Jika hujan maka terjadi banjir atau tanah subur. \\
    Banjir disebabkan oleh sampah dan selokan kecil. \\
    Tanah yang subur membuat pohon berbuah. \\ \\
    Kesimpulannya adalah $\dots$
    \begin{enumerate}
    \item Hujan menyebabkan sampah dan pohon berbuah.
    \item Hujan disebabkan sampah atau selokan kecil.
    \item \textbf{Banjir terjadi karena hujan atau sampah.}
    \item Banjir membuat pohon berbuah.
    \item Hujan membuat tanah subur dan banjir.
    \end{enumerate}
    \textbf{Penjelasan: } \\
    Misal $H$: Hujan, $B$: Banjir, $M$: Sampah. \\
    Premis 1: $H \implies (B \lor \dots)$ \textit{(``Hujan adalah salah satu sebab Banjir'')}. \\
    Premis 2: $(M \land \dots) \implies B$ \textit{(``Sampah adalah sebab lain dari Banjir'')}. \\
    Berdasarkan relasi kausalitas, kejadian Banjir ($B$) dapat terjadi akibat dipicu oleh Hujan ($H$) ATAU sampah ($M$). Secara logika: $B \impliedby (H \lor M)$.
  \item Mobil AC tidak terjadi embun kaca. \\
    Embun kaca terjadi jika suhu di dalam lebih panas daripada suhu di luar mobil. \\
    Suhu di luar mobil lebih dingin saat terjadi hujan. \\ \\
    Kesimpulannya adalah $\dots$
    \begin{enumerate}
    \item Mobil AC saat hujan terjadi embun kaca.
    \item Mobil non AC terjadi embun kaca saat tidak hujan.
    \item Mobil AC terjadi embun kaca saat suhu di dalam mobil lebih panas.
    \item \textbf{Mobil non AC tidak terjadi embun kaca saat suhu di luar mobil lebih panas.}
    \item Embun kaca tidak terjadi saat suhu di dalam mobil lebih panas.
    \end{enumerate}
    \textbf{Penjelasan: } \\
    Misal $P$: Suhu di dalam lebih panas dari luar. $E$: Terjadi embun. \\
    Premis 1: $P \implies E$ \textit{(``Jika dalam lebih panas, maka embun'')}. \\
    Berdasarkan aturan Kontraposisi: $\neg E \rightarrow \neg P$. \\
    Secara ekuivalen: $\neg P \rightarrow \neg E$. Artinya jika suhu luar lebih panas dari dalam ($\neg P$), maka pasti TIDAK terjadi embun kaca ($\neg  E$). Aturan alamiah ini berlaku general, termasuk untuk mobil non AC.
  \end{enumerate}
\end{multicols}
